%% ----------------------------------------------------------------------------
% CVG SA/MA thesis template
%
% Created 03/08/2024 by Tobias Fischer
%% ----------------------------------------------------------------------------

\chapter{Introduction}
\label{chap:intro}

\section{Motivation}

Many robots, vehicles, and wearable devices do not see an environment once --- they see it
repeatedly. A delivery robot drives the same warehouse floor every day; a fleet of survey
vehicles maps the same city over months; a person wearing an AR headset walks through the
same building on many separate occasions. Each of these passes through the environment ---
each \emph{session} --- is tracked independently by the device's own onboard visual or
visual-inertial odometry. Within a session, this odometry is locally accurate but
accumulates drift over time, and, critically, every session's trajectory is expressed in its
own local coordinate frame with no inherent relationship to any other session's frame: two
sessions that physically revisit the same corridor or intersection have no way, from
odometry alone, of knowing that they did.

Recovering that relationship --- detecting where sessions overlap and using those overlaps
to correct each session's drift and place every session into one shared, globally consistent
frame --- is the multi-session trajectory fusion problem. Solving it well is what turns a
pile of independent, drifting trajectory logs into a single coherent map of an environment,
and it is a prerequisite for any system that wants to accumulate knowledge about a place
across repeated visits rather than starting from zero every time.

\section{Technical Gap}

The classical way to detect that two sessions overlap is to also build and match dense 3D
scene geometry between them, as in multi-map SLAM systems such as ORB-SLAM3's Atlas
\cite{campos2021orbslam3}. Dense reconstruction gives strong evidence for a match, but it is
compute- and storage-heavy: keeping and matching dense maps across many long sessions scales
poorly, which limits deployments to a bounded number of sessions or a bounded environment
size.

A recent line of work sidesteps dense reconstruction entirely. Lim et al.'s 2GO
\cite{lim2025twogo}
represents each session only as a sparse, keyframed pose graph: cross-session overlap is
established with learned image retrieval and two-view geometry --- sparse feature matching
lifted to metric scale with a monocular depth prior, rather than any dense map --- and
fused directly into the pose graph. This is substantially lighter-weight than dense
reconstruction and is the closest prior work to this thesis.

2GO, however, is built around an offline, batch setting: every session's data is assumed
collected and available before fusion starts, and its formulation designates one session in
advance as the fixed reference frame the rest are fused into. Neither assumption holds for a
system meant to run continuously: sessions arrive one at a time, with no way to know in
advance which one --- if any single one at all --- should anchor the others. A genuinely
online multi-session system also has to decide, incrementally and without hindsight, how
much evidence is enough before it trusts that two sessions actually overlap; committing a
single spurious cross-session match into a shared pose graph can distort every trajectory
fused into it, not just the two sessions involved.

\section{This Work}

This thesis presents \textbf{Polygraph}, a multi-session trajectory fusion pipeline that
targets this online setting directly. Polygraph processes each session's odometry keyframes
incrementally, grouping them into local spatial clusters as they arrive; whenever a cluster
closes, it is matched against every other session's keyframes through learned retrieval, and
candidate cross-session matches are verified with the same sparse-matching-plus-monocular-depth
approach as 2GO. The key difference is architectural: Polygraph makes no assumption that all
sessions are known up front and designates no reference session in advance --- every session
starts as its own independent graph component, and cross-session edges merge components
together only once a two-tier consensus gate has corroborated them with independent evidence,
whether that session is being anchored into the graph for the first time or is already
connected. The shared pose graph is kept current online through a cheap incremental solve
after every closed cluster, and is periodically re-solved and pruned of outlier edges through
a robust batch optimization pass.

Concretely, this thesis contributes:
\begin{itemize}
  \item An online reformulation of sparse, retrieval-driven multi-session pose-graph fusion
    that removes the offline, pre-designated-reference-session assumption of the closest
    prior work, so that sessions can be processed incrementally as they are collected.
  \item An evaluation spanning KITTI; two indoor multi-session deployments built from the
    CroCoDL \cite{blum2025crocodl} and LaMAR \cite{sarlin2022lamar} benchmarks (captured on a
    HoloLens and a Boston Dynamics Spot respectively); and Oxford RobotCar with RTK ground
    truth, using a paired easy/hard session split --- sharing all but one session, with the
    hard split adding a night traversal --- to test both single- and multi-session performance
    under real day/night appearance change. None of these settings are covered by the closest
    prior work's own published evaluation.
  \item A direct, hardware-matched comparison against that prior work's published numbers,
    reporting honestly where Polygraph is faster and where an open gap in loop-closure recall
    remains, rather than only reporting favorable comparisons.
\end{itemize}

\section{Thesis Outline}

\mychapterref{chap:relatedwork} situates Polygraph against prior work in visual SLAM,
place recognition, learned two-view geometry, and multi-session mapping, with particular
attention to 2GO as the closest prior system. \mychapterref{chap:method} describes
Polygraph's pipeline in detail, stage by stage: keyframing, online clustering, retrieval-based
cluster matching, joint pose estimation, the two-tier consensus gate, and pose-graph
optimization. \mychapterref{chap:experiments} presents the experimental setup and results
across all evaluated datasets. \mychapterref{chap:discussion} discusses limitations,
including the open loop-closure-recall gap identified against 2GO, and design trade-offs made
along the way. \mychapterref{chap:conclusion} summarizes the thesis's contributions and
outlines directions for future work.

% REVIEW NOTE (not rendered): \mychapterref{} calls above assume \label{chap:relatedwork},
% \label{chap:method}, \label{chap:experiments}, \label{chap:discussion}, and
% \label{chap:conclusion} get added to the corresponding chapter headings in
% 02_relatedwork.tex .. 06_conclusion.tex (mirroring \label{chap:intro} above) -- currently
% only this chapter has one, so those refs will show as "??" until then.
%
% Also: RobotCar's results are being finalized (run overnight per our last conversation) --
% the contributions bullet above describes the *setup* (paired easy/hard split, day/night,
% single/multi-session), not results, which is safe regardless of how the numbers land. EuRoC
% is deliberately left out of this bullet -- its Atlas-comparison results (ORB-SLAM3 backbone)
% are done per docs/next_experiments.md item 5, but VINS-Fusion-on-EuRoC is still open and its
% overall place in the thesis narrative isn't settled yet; add it once that's decided.
